% Unit 099 terjemahan Bahasa Indonesia dari chapter7.tex baris 2036--2140.
% Sumber: Methods of Algebra, Volume 2, commit
% 9a5803ff77dd3257484cb177f851a73770a59dd3 (CC BY 4.0).
% stable-unit-id: o014.aljabr2.chapter7.identifying-module-categories
% Cakupan: Bagian 7.8 lengkap, ``Mengenali Kategori Modul'';
% berhenti tepat sebelum Bagian 7.9 pada baris 2141.

% segment-id: o014.li2.chapter7.identifying-module-categories.h001
\section{Mengenali Kategori Modul}\label{sec:module-cat}
\hypertarget{o014.aljabr2.chapter7.identifying-module-categories}{ }
% segment-id: o014.li2.chapter7.identifying-module-categories.p001
Dalam bagian ini, tetap dipilih gelanggang komutatif $\Bbbk$. Semua kategori
abelian dan funktor secara baku dianggap $\Bbbk$-linear.

% segment-id: o014.li2.chapter7.identifying-module-categories.p002
Kategori berbentuk $R\dcate{Mod}$ atau $\cated{Mod}R$ secara kolektif disebut
kategori modul, dengan $R$ suatu aljabar-$\Bbbk$. Dalam
\S\ref{sec:Morita}, kita secara khusus membahas kategori modul beserta
funktor di antaranya. Di sisi lain, wajar pula ditanyakan kategori mana
yang ekuivalen dengan kategori modul dan bagaimana ekuivalensinya diberikan
secara konkret. Bagian ini menyajikan beberapa hasil singkat dan praktis
tentang pertanyaan tersebut.

% segment-id: o014.li2.chapter7.identifying-module-categories.p003
Misalkan $\mathcal{A}$ kategori abelian. Pendekatan bagian ini ialah meninjau
funktor berbentuk $\Hom_{\mathcal{A}}(s,\mathord\cdot)$, yang secara alami
diangkat menjadi funktor $\mathcal{A}\to\cated{Mod}R$ dengan
$R:=\End_{\mathcal{A}}(s)$. Persoalannya ialah menentukan kapan
$\Hom_{\mathcal{A}}(s,\mathord\cdot)$ merupakan ekuivalensi.

\index{objek kompak}
% segment-id: o014.li2.chapter7.identifying-module-categories.p004
Kita ingat kembali Definisi~\ref{def:cpt-objects} tentang objek kompak.
Misalkan $\mathcal{A}$ mempunyai semua $\varinjlim$ kecil terfilter dan
$X\in\Obj(\mathcal{A})$. Objek $X$ disebut \emph{kompak} jika, untuk setiap
kategori kecil terfilter $I$ dan funktor $\alpha:I\to\mathcal{A}$, morfisme
kanonik berikut adalah isomorfisme:
% segment-id: o014.li2.chapter7.identifying-module-categories.q001
\[ \varinjlim_i \Hom\left( X, \alpha(i) \right) \to \Hom\left( X, \varinjlim \alpha \right). \]

% segment-id: o014.li2.chapter7.identifying-module-categories.le001
\begin{lemma}\label{prop:proj-compact}
	Misalkan $X$ objek projektif dalam kategori abelian kokomplet $\mathcal{A}$.
	Maka $X$ kompak jika dan hanya jika
	$\Hom_{\mathcal{A}}(X,\mathord\cdot)$ mempertahankan semua jumlah langsung
	kecil; yaitu, untuk sembarang himpunan kecil $I$ dan keluarga objek
	$(Y_i)_{i\in I}$, berlaku
	$\bigoplus_{i\in I}\Hom_{\mathcal{A}}(X,Y_i)\rightiso
	\Hom_{\mathcal{A}}(X,\bigoplus_{i\in I}Y_i)$. Bila kondisi tersebut
	berlaku, $\Hom_{\mathcal{A}}(X,\mathord\cdot)$ sebenarnya mempertahankan
	semua $\varinjlim$ kecil.
\end{lemma}
% segment-id: o014.li2.chapter7.identifying-module-categories.pf001
\begin{proof}
	Misalkan $X$ kompak. Jumlah langsung yang diindeks oleh himpunan kecil $I$
	adalah $\varinjlim$ terfilter dari jumlah langsung hingga (ambil semua
	subhimpunan hingga dari $I$, terurut oleh inklusi). Karena
	$\Hom_{\mathcal{A}}(X,\mathord\cdot)$ diketahui mempertahankan jumlah
	langsung hingga, funktor itu juga mempertahankan jumlah langsung kecil.

	Sebaliknya, syarat projektif sama artinya dengan
	$\Hom_{\mathcal{A}}(X,\mathord\cdot)$ mempertahankan kokernel, sedangkan
	$\varinjlim$ kecil umum dapat dikonstruksi dari jumlah langsung kecil dan
	kokernel. Jadi, $\Hom_{\mathcal{A}}(X,\mathord\cdot)$ mempertahankan semua
	$\varinjlim$ kecil.
\end{proof}

% segment-id: o014.li2.chapter7.identifying-module-categories.eg001
\begin{example}
	Misalkan $R$ aljabar-$\Bbbk$ dan $M$ modul-$R$ kanan projektif. Maka $M$
	kompak jika dan hanya jika $M$ dibangkitkan secara hingga.
	\begin{itemize}
		\item Keniscayaan: terdapat himpunan kecil $I$ dan homomorfisme
		$i:M\to R^{\oplus I}$ serta $p:R^{\oplus I}\to M$ dengan
		$pi=\identity_M$. Kekompakan mengakibatkan citra $i$ terkandung dalam
		$R^{\oplus I_0}$ untuk suatu subhimpunan hingga $I_0\subset I$. Ambil
		$p':=p|_{R^{\oplus I_0}}$; tetap berlaku $p'i=\identity_M$. Jadi, $M$
		dapat direalisasikan sebagai suku langsung modul bebas berperingkat
		hingga.
		\item Kecukupan: terapkan kriteria Lema~\ref{prop:proj-compact}. Setiap
		homomorfisme $M\to\bigoplus_{i\in I}Y_i$ niscaya mendarat dalam
		$\bigoplus_{i\in I_0}Y_i$ untuk suatu subhimpunan hingga
		$I_0\subset I$ (cukup memeriksa citra para pembangkit). Karena itu, $M$
		kompak.
	\end{itemize}
\end{example}

% segment-id: o014.li2.chapter7.identifying-module-categories.pr001
\begin{proposition}\label{prop:proj-compact-gen}
	Misalkan $\mathcal{A}$ kategori abelian kokomplet. Maka
	$s\in\Obj(\mathcal{A})$ adalah pembangkit projektif kompak jika dan hanya
	jika $\Hom_{\mathcal{A}}(s,\mathord\cdot):
	\mathcal{A}\to\Bbbk\dcate{Mod}$ merupakan funktor eksak setia yang
	mempertahankan semua $\varinjlim$ kecil.
\end{proposition}
% segment-id: o014.li2.chapter7.identifying-module-categories.pf002
\begin{proof}
	Menurut Proposisi~\ref{prop:injective-cogenerator}, menjadi pembangkit
	projektif ekuivalen dengan sifat eksak setia dari
	$\Hom_{\mathcal{A}}(s,\mathord\cdot)$. Menurut
	Lema~\ref{prop:proj-compact}, dengan asumsi projektif, kekompakan ekuivalen
	dengan fakta bahwa $\Hom_{\mathcal{A}}(s,\mathord\cdot)$ mempertahankan
	semua $\varinjlim$ kecil.
\end{proof}

% segment-id: o014.li2.chapter7.identifying-module-categories.th001
\begin{theorem}[P.\ Gabriel]\label{prop:identify-ModR}
	Misalkan $s\in\Obj(\mathcal{A})$ dan tetapkan
	$R:=\End_{\mathcal{A}}(s)$. Maka
	% segment-id: o014.li2.chapter7.identifying-module-categories.q002
	\[ \Hom_{\mathcal{A}}(s, \cdot): \mathcal{A} \to \cated{Mod}R \]
	merupakan ekuivalensi jika dan hanya jika $s$ adalah pembangkit projektif
	kompak bagi $\mathcal{A}$.
\end{theorem}
% segment-id: o014.li2.chapter7.identifying-module-categories.pf003
\begin{proof}
	Untuk menyingkat notasi, tuliskan
	$G:=\Hom_{\mathcal{A}}(s,\mathord\cdot):\mathcal{A}\to\cated{Mod}R$.
	Mula-mula, kita buktikan keniscayaannya. Karena $G$ adalah ekuivalensi dan
	$G(s)=R$, cukup memeriksa bahwa $R$ adalah pembangkit projektif kompak bagi
	$\cated{Mod}R$, yang langsung terlihat.

	Sekarang kita buktikan kecukupannya. Misalkan $s$ pembangkit projektif kompak.
	Proposisi~\ref{prop:proj-compact-gen} mengakibatkan $G$ adalah funktor eksak
	setia yang mempertahankan semua $\varinjlim$ kecil, karena komposisinya
	dengan funktor pelupa $\cated{Mod}R\to\Bbbk\dcate{Mod}$ memiliki sifat
	tersebut. Karena kesetiaan telah diketahui, cukup dibuktikan bahwa:
	\begin{compactitem}
		\item untuk semua $X,Y\in\Obj(\mathcal{A})$, pemetaan yang bersesuaian
		$\Hom_{\mathcal{A}}(X,Y)\to\Hom_R(GX,GY)$ surjektif;
		\item $G$ surjektif secara esensial.
	\end{compactitem}

	Pertama-tama, untuk setiap $X\in\Obj(\mathcal{A})$, definisikan morfisme
	$\phi_X:s^{\oplus\Hom_{\mathcal{A}}(s,X)}\to X$ yang pada suku langsung
	yang bersesuaian dengan $f:s\to X$ sama dengan $f$. Kita tunjukkan bahwa
	$\phi_X$ surjektif. Andaikan sebaliknya bahwa
	$\Coker(\phi_X)\neq0$. Sifat pembangkit memberikan morfisme tak nol
	$s\to\Coker(\phi_X)$ (Proposisi~\ref{prop:injective-cogenerator}), sedangkan
	sifat projektif membuat morfisme tersebut terfaktorkan sebagai
	$s\xrightarrow{g}X\to\Coker(\phi_X)$. Akibatnya, komposisi
	$s^{\oplus\Hom_{\mathcal{A}}(s,X)}\xrightarrow{\phi_X}X
	\to\Coker(\phi_X)$ tidak nol pada suku langsung yang bersesuaian dengan
	$g$, suatu kontradiksi.

	Konstruksi yang sama dapat diterapkan pada kernel $\Ker(\phi_X)$. Jadi,
	setiap $X$ dapat ditempatkan dalam barisan eksak
	% segment-id: o014.li2.chapter7.identifying-module-categories.q003
	\[ s^{\oplus J} \to s^{\oplus I} \to X \to 0, \quad I = I_X, \; J = J_X: \text{himpunan kecil}. \]
	Karena $G$ mempertahankan semua $\varinjlim$ kecil, diperoleh barisan eksak
	modul-$R$ kanan
	% segment-id: o014.li2.chapter7.identifying-module-categories.q004
	\[ R^{\oplus J} \to R^{\oplus I} \to GX \to 0. \]
	Morfisme pertama dalam kedua barisan dapat dipandang sebagai matriks
	$I\times J$ di atas $R$ yang beraksi melalui perkalian kiri; keduanya
	sebenarnya adalah matriks yang sama.

	Sekarang, tinjau homomorfisme modul-$R$ kanan $f:GX\to GY$. Pilih untuk
	$X$ dan $Y$ barisan eksak berbentuk seperti di atas. Mudah dilihat bahwa
	$f$ dapat diangkat menjadi diagram komutatif dengan baris-baris eksak
	% segment-id: o014.li2.chapter7.identifying-module-categories.q005
	\[\begin{tikzcd}
		R^{\oplus J_X} \arrow[r] \arrow[d] & R^{\oplus I_X} \arrow[d] \arrow[r] & GX \arrow[d, "f"] \arrow[r] & 0 \\
		R^{\oplus J_Y} \arrow[r] & R^{\oplus I_Y} \arrow[r] & GY \arrow[r] & 0.
	\end{tikzcd}\]
	Namun, keempat anak panah pada persegi kiri dapat dinyatakan sebagai
	matriks di atas $R$. Keempat matriks itu memberikan bagian bergaris penuh
	dari diagram komutatif dengan baris-baris eksak
	% segment-id: o014.li2.chapter7.identifying-module-categories.q006
	\[\begin{tikzcd}
		s^{\oplus J_X} \arrow[r] \arrow[d] & s^{\oplus I_X} \arrow[d] \arrow[r] & X \arrow[dashed, d] \arrow[r] & 0 \\
		s^{\oplus J_Y} \arrow[r] & s^{\oplus I_Y} \arrow[r] & Y \arrow[r] & 0,
	\end{tikzcd}\]
	dan bagian bergaris putus-putus ditentukan secara unik oleh funktorialitas
	kokernel. Karena $G$ eksak, niscaya $G(X\dashrightarrow Y)=f$.

	Teknik serupa menunjukkan bahwa $G$ surjektif secara esensial. Untuk modul-$R$
	kanan $M$ yang diberikan, terdapat barisan eksak
	% segment-id: o014.li2.chapter7.identifying-module-categories.q007
	\[ R^{\oplus L} \to R^{\oplus K} \to M \to 0, \]
	dengan morfisme pertamanya dipandang sebagai matriks $K\times L$ di atas
	$R$. Matriks yang sama memberikan barisan eksak kokernel dalam
	$\mathcal{A}$,
	% segment-id: o014.li2.chapter7.identifying-module-categories.q008
	\[ s^{\oplus L} \to s^{\oplus K} \to X \to 0. \]
	Karena $G$ mempertahankan semua $\varinjlim$ kecil, niscaya
	$GX\simeq M$. Selesailah bukti.
\end{proof}

% segment-id: o014.li2.chapter7.identifying-module-categories.p005
Persoalan terkait lainnya ialah mengenali kategori modul yang dibangkitkan
secara hingga. Tuliskan $\catesubd{Mod}{\mathrm{fg}}R$ untuk kategori modul-
$R$ kanan yang dibangkitkan secara hingga. Jika $R$ gelanggang Noether kanan,
$\catesubd{Mod}{\mathrm{fg}}R$ tetap merupakan kategori abelian; lihat
\cite[Proposisi 6.10.3, Lema 6.10.4]{Li1}.
\index[sym1]{Modfg@$\cate{Mod}_{\mathrm{fg}}$}

% segment-id: o014.li2.chapter7.identifying-module-categories.th002
\begin{theorem}\label{prop:identify-ModR-fg}
	Misalkan $\mathcal{A}$ kategori abelian dan $s\in\Obj(\mathcal{A})$. Andaikan:
	\begin{compactitem}
		\item setiap $X\in\Obj(\mathcal{A})$ adalah objek Noether
		(Definisi~\ref{def:finite-length-object});
		\item $R:=\End_{\mathcal{A}}(s)$ adalah gelanggang Noether kanan;
		\item $s$ adalah pembangkit projektif bagi $\mathcal{A}$.
	\end{compactitem}
	Maka funktor $\Hom_{\mathcal{A}}(s,\mathord\cdot):
	\mathcal{A}\to\cated{Mod}R$ terfaktorkan sebagai
	% segment-id: o014.li2.chapter7.identifying-module-categories.q009
	\[ \mathcal{A} \xrightarrow{\text{ekuivalensi}} \catesubd{Mod}{\mathrm{fg}}R \subset \cated{Mod}R. \]
\end{theorem}
% segment-id: o014.li2.chapter7.identifying-module-categories.pf004
\begin{proof}
	Pertama-tama, tunjukkan bahwa setiap $X\in\Obj(\mathcal{A})$ dapat
	ditempatkan dalam barisan eksak
	% segment-id: o014.li2.chapter7.identifying-module-categories.q010
	\begin{equation}\label{eqn:identify-ModR-fg-aux}
		s^{\oplus m} \to s^{\oplus n} \to X \to 0, \quad n, m \in \Z_{\geq 0}.
	\end{equation}

	Untuk itu, karena $s$ adalah pembangkit, bagi setiap $X\neq0$ terdapat
	morfisme tak nol $\phi_1:s\to X$. Jika $\phi_1$ tidak surjektif, terdapat
	morfisme tak nol $s\to\Coker(\phi_1)$; sifat projektif $s$ memastikan bahwa
	morfisme ini terfaktorkan sebagai
	$s\xrightarrow{\psi_1}X\to\Coker(\phi_1)$. Dengan demikian, diperoleh
	$\phi_2:=(\phi_1,\psi_1):s^{\oplus2}\to X$ dengan
	$\Image(\phi_2)\supsetneq\Image(\phi_1)$. Jika $\phi_2$ belum surjektif,
	ulangi prosedur yang sama. Karena $X$ objek Noether, prosedur ini harus
	berhenti setelah sejumlah langkah berhingga dan menghasilkan morfisme
	surjektif $\phi_X:s^{\oplus n}\twoheadrightarrow X$. Terapkan prosedur yang
	sama pada $\Ker(\phi_X)$ untuk memperoleh
	\eqref{eqn:identify-ModR-fg-aux}.

	Ingat bahwa $s$ adalah pembangkit projektif jika dan hanya jika
	$G:=\Hom_{\mathcal{A}}(s,\mathord\cdot)$ eksak setia. Menerapkan $G$ pada
	\eqref{eqn:identify-ModR-fg-aux} memberikan barisan eksak
	$R^{\oplus m}\to R^{\oplus n}\to GX\to0$, yang menunjukkan bahwa citra
	$G$ terletak dalam $\catesubd{Mod}{\mathrm{fg}}R$.

	Di sisi lain, karena $R$ gelanggang Noether kanan, setiap
	$M\in\Obj(\catesubd{Mod}{\mathrm{fg}}R)$ juga dapat ditempatkan dalam
	barisan eksak
	% segment-id: o014.li2.chapter7.identifying-module-categories.q011
	\[ R^{\oplus q} \to R^{\oplus p} \to M \to 0, \quad p, q \in \Z_{\geq 0}. \]

	Selanjutnya, gagasan pembuktiannya sama dengan bukti kecukupan dalam
	Teorema~\ref{prop:identify-ModR}. Dalam argumen tersebut, kekompakan $s$
	hanya digunakan untuk menangani jumlah langsung tak hingga; di sini hanya
	terlibat jumlah langsung hingga, yang dipertahankan oleh setiap funktor
	aditif. Selebihnya, argumennya persis sama dan tidak diulangi.
\end{proof}
